\section{An Entity Is a Type With a Key}
\label{sec:ch05-an-entity-is-a-type-with-a-key}

\index{entity!definition|(}%
Section~\ref{sec:ch05-the-join-you-cannot-see} left three integers stranded on
the wrong side of a boundary. An entity is the schema construct that gets them
across, and Apollo defines it in terms a reader of
chapter~\ref{ch:data-without-n-plus-1} will recognize~\autocite{apolloentities}:

\begin{quote}
Entities are objects that can be fetched with one or more unique key fields.
Like a row in a database table, an entity contains fields of various types, and
it can be uniquely identified by a key field or set of fields.
\end{quote}

\index{entity!requires a key}%
The clause about key fields is the one that does the work, and the same page
makes it a requirement rather than a description: \enquote{Because an entity is
keyed, an entity type's definition must have a \code{@key}
directive}~\autocite{apolloentities}. A type with no key is an ordinary object
type. It can appear in one subgraph's schema and be perfectly useful there, and
no second subgraph can say anything about it.

\begin{graphqlsdl}
directive @key(fields: FieldSet!, resolvable: Boolean = true)
  repeatable on OBJECT | INTERFACE
\end{graphqlsdl}

\index{FieldSet@\code{FieldSet}!is a selection set as a string}%
\code{FieldSet} is a scalar and the specification's description of it is the
best two lines in the document: \enquote{Grammatically, a \code{FieldSet} is a
selection set minus the outermost curly
braces}~\autocite{apollosubgraphspec}. So the key is written the way you would
write the query that fetches it. \code{"id"} is a key. So is
\code{"username region"}, two fields that are unique together and neither of
which is unique alone. So is \code{"name organization \{ id \}"}, which reaches
into a nested object~\autocite{apollodirectives}. A type may carry more than
one \code{@key}, which says there are two different ways to name the same
object, and the directive is \code{repeatable} for exactly that reason.

\subsection{What the Key Costs You}

\index{key!must be a field}%
Read the definition again for what it forces. A key names fields, and a field
is a thing in a schema. So every subgraph that declares the type has to declare
the key fields on it, in public, whether or not that subgraph has any use for
them.

\index{Speaker@\code{Speaker}!appears in a subgraph that knows nothing about speakers}%
For the conference example that means the Sessions subgraph, which owns no
speaker data and never will, declares a \code{Speaker} type anyway:

\begin{graphqlsdl}
type Speaker @key(fields: "id") {
  id: ID!
}
\end{graphqlsdl}

\index{stub!a type declared for its key alone}%
One field, and it is there so that the value can be written down. This is the
thing section~\ref{sec:ch05-the-join-you-cannot-see} predicted: the key stops
being an integer in a call stack and becomes a field with a name, a type and a
promise attached. Chapter~\ref{ch:data-without-n-plus-1} took a foreign key out
of the schema, and the model has just put one back, in a different shape and for
a reason that did not exist then.

\index{key!both sides must produce the same value}%
And both subgraphs have to produce the same string for the same speaker. Not an
equivalent one, not one that maps to it: the same bytes. The router matches on
that value and has no other way to know that Sessions' speaker one and Speakers'
speaker one are the same person. Nothing in the schema checks this and no
composition error catches it, because both sides declare \code{id: ID!} and
both are telling the truth about the type. Chapter~\ref{ch:representation-order}
is about a different way to get the wrong entity back, and this is the way that
is entirely yours.

\index{global object identification!is why the key already exists}%
The good news arrives from a chapter that was not thinking about federation at
all. Chapter~\ref{ch:schema-design} gave every entity a global node id, derived
from the type name and the primary key by a rule both services can run.
\code{U3BlYWtlcjox} is \code{Speaker:1} in every process that computes it the
same way. The book already has its keys.

\subsection{The Route the Key Opens}

\index{entities@\code{\_entities}!takes representations}%
A key is half a mechanism. It lets the router write down which object it means;
it does not yet let anyone turn that back into an object. The other half is the
entity route, and it is one root field:

\begin{graphqlsdl}
extend type Query {
  _entities(representations: [_Any!]!): [_Entity]!
  _service: _Service!
}
\end{graphqlsdl}

\index{Any@\code{\_Any}!is a JSON object}%
\code{\_Any} is deliberately shapeless. The specification says each item is
\enquote{serialized as a generic JSON object, which enables the graph router to
include representations of different entities in the same query, all of which
can have a different shape}~\autocite{apollosubgraphspec}. A representation is
the type name plus the key fields, and nothing else:

\begin{minted}{json}
{"__typename": "Speaker", "id": "U3BlYWtlcjox"}
\end{minted}

\index{Entity@\code{\_Entity}!is generated from the keys}%
\code{\_Entity} is a union the library generates for you, holding every type in
the schema that carries a resolvable key. It has to be a union because one call
can carry representations of several types at once, which is also why the router
can collect a whole level of a query into a single request.

\index{reference resolver!is the inverse of the key}%
What you write is one function per entity, and its shape is the inverse of the
key: the key turns an object into a set of values, and this turns that set of
values back into the object. Hot Chocolate calls it a reference resolver and
chapter~\ref{ch:first-subgraph} is where you write the first one. Note now what
it is not. It is not a resolver for a field, it does not know which client
asked, and it does not know which field of which query caused it to run. It is
handed a key and asked for the object.

\index{entities@\code{\_entities}!answers in order}%
One rule about it belongs here rather than there, because it is a property of
the contract rather than of any implementation. The specification requires that
\code{\_entities} \enquote{return a list of entity objects that correspond to
the provided representations, in the exact same order}, and permits nulls in
that list \enquote{if no entity exists for a provided
representation}~\autocite{apollosubgraphspec}. Position is the only thing
connecting an answer to its question. There is no identifier in the reply, no
echo of the representation, nothing to match on. Get the order wrong and every
answer is still well typed, still non-null where the schema promised, and
attached to the wrong object. Chapter~\ref{ch:representation-order} is that
chapter.

\index{key@\code{"@key}!resolvable argument}%
Which leaves \code{resolvable}. Set it to false and you are declaring a key
without promising to answer for it: the specification's words are that the
subgraph \enquote{doesn't define a reference resolver for this entity}, so the
router's plans cannot route into this subgraph to fetch the type's
fields~\autocite{apollodirectives}. That is the right setting for a type you
only ever point at, and the wrong setting for a type you own. It is also the
one argument on \code{@key} whose default, \code{true}, is a promise you make
by saying nothing.
\index{entity!definition|)}%
